\chapter{Theoretischer Hintergrund}
TODO: Einleitung für das Kapitel schreiben

\section{Definition: Webanwendung}
Eine Webanwendung (auch \emph{Web-App} genannt) ist eine Anwendung, die vollständig in einem Browser ausführbar ist. Da sie (fast) vollständig auf einem Webserver ausgeführt wird, ist das zugrunde liegende Übertragungsprotokoll \emph{HTTP}.\par

Web-Apps bieten den Vorteil, dass sie in jedem modernen Browser ausgeführt werden können, ohne, dass spezielle Programmiersprachen erlernt und angewendet werden müssen. Interessant sind diese Anwendungen, da sie auch auf Smartphones und Tablets ausgeführt werden können ohne sie vollständig auf die jeweilige Plattform portieren zu müssen.

\section{HTML5}
\begin{figure}[!ht]
	\centering
	\includegraphics[width=3cm]{fig/html5_logo}
	\caption{{HTML5-Logo des W3C}\cite{html5:logo}}
\end{figure}